\documentclass[12pt,aspectratio=1610,professionalfont,serif,t]{beamer}
\usepackage[osf,p]{newpxtext}
\usepackage{newpxmath}
\usepackage{fontawesome5}

\usepackage{fancyvrb}
\fvset{%frame=single,
commandchars=\\\{\},
%framesep=1mm,
fontfamily=helvetica,
fontsize=\normalsize
}
%\usepackage[bitstream-charter]{mathdesign}
\usepackage{listings}
\lstset{commentstyle=\color{orange}, keywordstyle=\color{yellow} }


\usetheme{Boadilla}
\useoutertheme{split}
\usecolortheme{albatross}


\setbeamertemplate{blocks}[rounded][shadow=false]
\setbeamersize{text margin left=\dimexpr 1cm+1ex\relax, text margin right=\dimexpr 1cm+1ex\relax}
\addtobeamertemplate{headline}{}{\usebeamerfont{frametitle}\vspace*{1ex}}
\setbeamertemplate{navigation symbols}{}
\setbeamertemplate{page number in head/foot}{}
%\beamerdefaultoverlayspecification{<+->}


\setbeamercolor*{structure}{fg=green!60!black,bg=black}
\setbeamercolor*{normal text}{fg=green!60!black,bg=black}
\setbeamercolor{palette primary}{use={structure,normal text},fg=green,bg=structure.bg}
\setbeamercolor{palette secondary}{use={structure,normal text},fg=structure.fg,bg=structure.bg}
\setbeamercolor{palette tertiary}{use={structure,normal text},fg=structure.fg,bg=structure.bg}
\setbeamercolor{palette quaternary}{use={structure,normal text},fg=white,bg=green!30!black}
\setbeamercolor{palette sidebar quaternary}{use={structure,normal text},fg=white,bg=green!30!black}
\setbeamercolor{title in head/foot}{use={structure,normal text},fg=white,bg=green!30!black}

\setbeamercolor*{example text}{use={structure},fg=green!60!black}
\setbeamercolor*{block body}{use={structure},bg=structure.bg}
\setbeamercolor*{block body alerted}{use={structure},bg=structure.bg}
\setbeamercolor*{block body example}{use={structure},bg=structure.bg}
\setbeamercolor*{block title}{parent=structure,bg=structure.bg}
\setbeamercolor*{block title alerted}{use={structure,alerted text},fg=alerted text.fg,bg=structure.bg}
\setbeamertemplate{navigation symbols}{}
\setbeamertemplate{items}[square]
\setbeamercolor{item projected}{fg=white}
\setbeamercolor*{normal text}{use={structure},fg=white!90!blue,bg=structure.bg}
\setbeamercolor*{separation line}{}
\setbeamercolor*{fine separation line}{}
\setbeamercolor{alerted text}{use={structure},fg=structure.fg}

\usepackage[italian]{babel}
\usepackage[utf8]{inputenc}
\usepackage{pgf}
\usepackage{verbatim}
\usepackage{inconsolata}
\usepackage{listings}
\lstset{language=C, frame=single,
	basicstyle=\ttfamily,
	numbers=left, numberstyle=\tiny\color{gray},
	numbersep=5pt, fancyvrb=true
}

\usepackage[noend]{algorithm2e}
\usepackage{pgf}
\usepackage{tikz}
\usepackage{graphicx}
\usetikzlibrary{%
	arrows,
	arrows.meta,
	positioning,
	calc,
	backgrounds,
	chains,
	matrix,
	patterns,
	automata,
	fit,
	graphs,
	decorations,
	decorations.pathmorphing,
	decorations.pathreplacing,
	decorations.markings,
}

\usepgflibrary{shapes,shapes.geometric}


\usepackage{url}


\usepackage{xmpmulti}
\usepackage{transparent}
\graphicspath{{img/}}

\newcommand{\zero}{\ensuremath{\mathtt{0}}\xspace}
\newcommand{\uno}{\ensuremath{\mathtt{1}}\xspace}
\newcommand{\due}{\ensuremath{\mathtt{2}}\xspace}
\newcommand{\tre}{\ensuremath{\mathtt{3}}\xspace}
\newcommand{\kleene}{\ensuremath{{}^{\ast}}\xspace}
\xspaceaddexceptions{\zero \uno \due \tre \kleene}


\newcommand{\simplebox}[1]{
      \begin{tikzpicture}
              \node[text centered, rectangle, draw=black, thick, text width=\linewidth,align=center, rounded corners, minimum height=1em] (C) {#1};
      \end{tikzpicture}
}
%
% In questo file si mettono i comandi utili solo per questo insegnamento.
%
\author{Gianluca Della Vedova}
\title{Linguaggi e Computabilità}
\institute{Univ. Milano--Bicocca\\
	\texttt{https://www.unimib.it/gianluca-della-vedova}}
\date{\today}



\begin{document}


\begin{frame}[fragile]
        \frametitle{Linguaggi e Computabilità: Lezione 17}
\begin{itemize}
\item
      Automi a pila: definizione
\item
      Rappresentazione grafica
\item
      Configurazione
\item
      Accettazione per stato finale e per pila vuota
\end{itemize}
\end{frame}



% \begin{frame}
%         \frametitle{Notazione}
%         Un arco da $p$ a $q$ con etichetta
%         \[
%                 a, X / \gamma
%         \]
%         significa $(q, \gamma) \in \delta(p, a, X)$:
%         \begin{itemize}
%                 \item si legge $a \in \Sigma \cup \{\varepsilon\}$ dall'input
%                 \item si toglie $X$ dalla cima della pila
%                 \item si mette $\gamma \in \Gamma\kleene$ sulla pila (il primo simbolo di $\gamma$ va in cima)
%         \end{itemize}

%         \bigskip
%         Configurazione (descrizione istantanea): $(q, w, \alpha)$, con $q$ stato corrente, $w$ input da
%         leggere, $\alpha$ contenuto della pila (cima a sinistra).
% \end{frame}



\begin{frame}
        \frametitle{Esempio: accettazione per stato finale}
        $L = \{ \zero^{n} \uno^{n} \; \vert \; n \ge 0 \}$

        \smallskip
        \begin{center}
                \begin{tikzpicture}[enfa, node distance=4cm]
                        \node[state, initial]          (q0)                {$q_0$};
                        \node[state]                   (q1) [right=of q0]  {$q_1$};
                        \node[state, accepting]        (q2) [right=of q1]  {$q_2$};

                        \path (q0) edge[loop above] node[align=center]
                                        {$\zero, Z_0 / \zero Z_0$ \\ $\zero, \zero / \zero\zero$} ()
                                   edge node {$\uno, \zero / \varepsilon$} (q1)
                                   edge[bend right=30] node[below] {$\varepsilon, Z_0 / Z_0$} (q2)
                              (q1) edge[loop above] node {$\uno, \zero / \varepsilon$} ()
                                   edge node {$\varepsilon, Z_0 / Z_0$} (q2);
                \end{tikzpicture}
        \end{center}
\end{frame}



\begin{frame}
        \frametitle{Esempio: computazione}
        % Parola $\zero\zero\uno\uno \in L$:
        % \[
        %         \begin{array}{l}
        %                 (q_0, \zero\zero\uno\uno, Z_0) \vdash
        %                 (q_0, \zero\uno\uno, \zero Z_0) \vdash
        %                 (q_0, \uno\uno, \zero\zero Z_0) \vdash \\
        %                 (q_1, \uno, \zero Z_0) \vdash
        %                 (q_1, \varepsilon, Z_0) \vdash
        %                 (q_2, \varepsilon, Z_0)
        %         \end{array}
        % \]
        % Input esaurito in uno stato finale: la parola è accettata.

        % \bigskip
        % Parola $\zero\zero\uno \notin L$:
        % \[
        %         (q_0, \zero\zero\uno, Z_0) \vdash
        %         (q_0, \zero\uno, \zero Z_0) \vdash
        %         (q_0, \uno, \zero\zero Z_0) \vdash
        %         (q_1, \varepsilon, \zero Z_0)
        % \]
        % Nessuna mossa possibile e $q_1$ non è finale.

        % \bigskip
        % Parola $\zero\uno\uno \notin L$: si arriva in $(q_1, \uno, Z_0)$, si può solo andare in
        % $(q_2, \uno, Z_0)$, ma l'input non è esaurito.
\end{frame}




\begin{frame}
        \frametitle{Esempio: accettazione per pila vuota}
        $L = \{ \zero^{n} \uno^{n} \; \vert \; n \ge 0 \}$

        \smallskip
        \begin{center}
                \begin{tikzpicture}[enfa, node distance=4cm]
                        \node[state, initial]          (q0)                {$q_0$};
                        \node[state]                   (q1) [right=of q0]  {$q_1$};
                        \node[state]        (q2) [right=of q1]  {$q_2$};

                        \path (q0) edge[loop above] node[align=center]
                                        {$\zero, Z_0 / \zero Z_0$ \\ $\zero, \zero / \zero\zero$} ()
                                   edge node {$\uno, \zero / \varepsilon$} (q1)
                                   edge[bend right=30] node[below] {$\varepsilon, Z_0 / \varepsilon$} (q2)
                              (q1) edge[loop above] node {$\uno, \zero / \varepsilon$} ()
                                   edge node {$\varepsilon, Z_0 / \varepsilon$} (q2);
                \end{tikzpicture}
        \end{center}
        \end{frame}

\begin{frame}
        \frametitle{Esercizio 1: accettazione per pila vuota}
Scrivere un PDA che accetta per pila vuota il linguaggio        $L = \{ w w^{R} \; \vert \; w \in \Sigma_{B}\kleene \}$
        % \begin{center}
        %         \begin{tikzpicture}[enfa, node distance=5cm]
        %                 \node[state, initial]  (q0)                {$q_0$};
        %                 \node[state]           (q1) [right=of q0]  {$q_1$};
        %                 \path (q0) edge[loop above] node[align=center]
        %                                 {$\zero, X / \zero X$ \\ $\uno, X / \uno X$} ()
        %                            edge node {$\varepsilon, X / X$} (q1)
        %                       (q1) edge[loop above] node[align=center]
        %                                 {$\zero, \zero / \varepsilon$ \\ $\uno, \uno / \varepsilon$ \\
        %                                  $\varepsilon, Z_0 / \varepsilon$} ();
        %         \end{tikzpicture}
        %         {\small($X$ indica un qualsiasi simbolo di $\Gamma = \{\zero, \uno, Z_0\}$)}
        % \end{center}
\end{frame}




\begin{frame}
        \frametitle{Esercizio 2}
        Scrivere un PDA che riconosce il linguaggio
        \[
                L = \{ \zero^{n} \uno^{m} \; \vert \; n \neq m, n\ge 0, m\ge 0 \}
        \]
        sull'alfabeto $\Sigma_{B}$.
% Soluzione (accettazione per stato finale), Gamma = {0, Z0}.
% Stati: q0 (lettura degli 0), q1 (lettura degli 1, pila con 0),
% q2 finale (n > m), q3 finale (n < m).
% q0: 0, Z0 / 0 Z0;  0, 0 / 00            (si impilano gli 0)
% q0 -> q1: 1, 0 / epsilon
% q1: 1, 0 / epsilon                        (ogni 1 toglie uno 0)
% q0 -> q2: epsilon, 0 / 0                  (n > 0, m = 0)
% q1 -> q2: epsilon, 0 / 0                  (input finito con 0 sulla pila: n > m)
% q0 -> q3: 1, Z0 / Z0                      (n = 0, m > 0)
% q1 -> q3: 1, Z0 / Z0                      (si legge un 1 a pila "vuota": n < m)
% q3: 1, Z0 / Z0                            (si leggono gli 1 restanti)
% q2 non ha transizioni uscenti: accetta solo se l'input è esaurito.
% Se n = m si termina in q0 (n = 0) o in q1 con Z0 in cima: stati non finali.
% Dopo un 1 non si può più leggere uno 0 (q1, q2, q3 non hanno transizioni con 0).
\end{frame}


\begin{frame}
        \frametitle{Esercizio 3}
        Scrivere un PDA che riconosce il linguaggio
        \[
                L = \{ \zero^{i} \uno^{j} \due^{k} \; \vert \; j = i + k, i>0, j>0 \}
        \]
        sull'alfabeto $\{\zero, \uno, \due\}$.
% Soluzione (accettazione per stato finale), Gamma = {0, 1, Z0}.
% Una parola di L si scompone in modo unico come 0^i 1^i 1^k 2^k:
% i primi i simboli 1 cancellano gli 0, gli altri k vengono impilati
% e poi cancellati dai 2.
% q0: 0, Z0 / 0 Z0;  0, 0 / 00             (si impilano gli 0)
% q0 -> q1: 1, 0 / epsilon                  (garantisce i > 0)
% q1: 1, 0 / epsilon                        (gli 1 cancellano gli 0)
% q1 -> q2: 1, Z0 / 1 Z0                    (finiti gli 0, si impilano gli 1)
% q2: 1, 1 / 11
% q2 -> q3: 2, 1 / epsilon
% q3: 2, 1 / epsilon                        (i 2 cancellano gli 1)
% q1 -> q4: epsilon, Z0 / Z0                (caso k = 0)
% q3 -> q4: epsilon, Z0 / Z0
% q4 unico stato finale, senza transizioni uscenti.
% Variante per pila vuota: q1 e q3 hanno epsilon, Z0 / epsilon al posto
% delle transizioni verso q4.
\end{frame}


\begin{frame}
        \frametitle{Esercizio 4}
        Scrivere un PDA che riconosce il linguaggio
        \[
                L = \{ \zero^{n} \uno^{m} \; \vert \; n \le m \le 2n, n>0, m>0 \}
        \]
        sull'alfabeto $\Sigma_{B}$.

        \bigskip

% Soluzione (accettazione per pila vuota), Gamma = {X, Z0}.
% Per ogni 0 letto si impilano, in modo non deterministico, una oppure due X;
% ogni 1 letto toglie una X.
% q0: 0, Z0 / X Z0;  0, Z0 / XX Z0;  0, X / XX;  0, X / XXX
% q0 -> q1: 1, X / epsilon                  (garantisce n > 0 e m > 0)
% q1: 1, X / epsilon
% q1: epsilon, Z0 / epsilon                 (si svuota la pila)
% Dopo aver letto 0^n la pila contiene tra n e 2n simboli X (ogni valore
% intermedio è possibile), quindi la pila si svuota leggendo esattamente
% m simboli 1 se e solo se n <= m <= 2n.
% Variante per stato finale: q1 -> q2 con epsilon, Z0 / Z0, q2 finale.
\end{frame}



\end{document}

%%% Local Variables:
%%% mode: LaTeX
%%% TeX-engine: luatex
%%% TeX-master: t
%%% End:
